\chapter{LaTeX-орієнтоване управління проєктами}

\section{Концепція продукту «ERP/1: Проєкти»}
У великих інженерних проєктах, державних замовленнях та системах подвійного призначення розробка та погодження технічної документації (технічних вимог --- ТВ, технічних завдань --- ТЗ, техноробочих проєктів --- ТРП) є невід'ємною частиною розробки програмного забезпечення.

Популярні інструменти управління проєктами, такі як Atlassian Jira (відстеження задач) та Confluence (база знань), використовують спрощені текстові розмітки (Markdown або Visual Rich Text), які не здатні забезпечити видавничу якість оформлення складних інженерних креслень, блок-схем, специфікацій сокетів та математичних формул.

Модуль «ERP/1: Проєкти» побудований як інтегрована заміна Jira та Confluence, що має такі ключові відмінності:
\begin{itemize}
    \item \textbf{Процесне управління (Jira-аналог)} --- робота з Agile-дошками (Scrum, Kanban) та гнучке налаштування статусів задач за допомогою процесного рушія BPMN (BPE).
    \item \textbf{LaTeX як єдиний стандарт розмітки (Confluence-аналог)} --- опис вимог до задач та вікі-статей бази знань виконується безпосередньо у форматі LaTeX.
    \item \textbf{Локальне розгортання} --- повна автономність та відсутність зовнішніх хмарних залежностей.
\end{itemize}

\section{Серверна інтеграція та PDF-компіляція}
База знань системи зберігає вихідний код LaTeX-документів у KVS/Mnesia. При кожному оновленні сторінки або за запитом користувача, фоновий Erlang-процес ініціює компіляцію документа за допомогою локального пакету TeX Live (\texttt{xelatex} або \texttt{pdflatex}).

Скомпільований PDF-документ кешується та відображається безпосередньо в інтерфейсі користувача за допомогою веб-переглядача PDF. Це дозволяє розробникам та технічним письменникам:
\begin{enumerate}
    \item Використовувати пакет \texttt{amsmath} для складних математичних рівнянь.
    \item Малювати точні векторні діаграми архітектур та мереж зв'язку за допомогою пакетів \texttt{tikz} та \texttt{pgf}.
    \item Автоматично генерувати готову технічну документацію за стандартами ДСТУ та ISO/IEC 12207 безпосередньо з тикетів та статей бази знань.
\end{enumerate}

\section{Переваги інтегрованої LaTeX-документації}
Впровадження LaTeX-орієнтованого управління проєктами дозволяє:
\begin{itemize}
    \item \textbf{Дотримуватись єдиного джерела істини (Single Source of Truth)}: технічні специфікації, що описують код, зберігаються безпосередньо поруч із кодом у тих самих репозиторіях і мають ідентичний вигляд як у системі ведення завдань, так і в друкованих звітах.
    \item \textbf{Гарантувати відсутність залежностей}: система не потребує сторонніх хмарних конвертерів чи важких JS-бібліотек для рендерингу складних математичних формул на клієнтських пристроях.
    \item \textbf{Спрощувати проходження сертифікації}: згенеровані PDF-файли ТВ та ТЗ повністю відповідають державним вимогам до паперового діловодства, що полегшує здачу робіт замовнику.
\end{itemize}

\section{Тестові питання та завдання}
\begin{enumerate}
    \item У чому полягають обмеження розмітки Markdown у системах управління великими інженерними проектами?
    \item Опишіть схему взаємодії фонового Erlang-процесу з локальним TeX Live для рендерингу документів бази знань.
    \item Які переваги дає використання векторного пакету TikZ для документування архітектури ПЗ?
    \item Як інтегрований генератор LaTeX-звітів допомагає при проходженні державної сертифікації програмних рішень?
    \item \textbf{Практичне завдання}: Напишіть простий шаблон LaTeX-документа для оформлення сторінки технічного завдання, що містить таблицю з переліком вимог до продуктивності системи.
\end{enumerate}
